\documentclass[10pt]{article}

\usepackage[utf8]{inputenc}

\usepackage[T1]{fontenc}

\usepackage{amsmath}

\usepackage{amsfonts}

\usepackage{amssymb}

\usepackage[version=4]{mhchem}

\usepackage{stmaryrd}

\usepackage{mathrsfs}



\begin{document}

(наз. соответственно о п е р а т о р а м и г р а н е й и опе раторами вы рождения). В случае, когда $\mathscr{C}$ является категориеї структуризованных множеств, точки множества $X_{n}$ наз. обычно $n$-м е pны м и с и мп л е к с а м и С. о. $X$. Отображения $\delta_{i}$ и $\sigma_{i}$ удовлетворяют соотношениям



\[

\begin{aligned}

\delta_{j} \delta_{i} & =\delta_{i} \delta_{j-1}, \text { если } i<j, \\

\sigma_{j} \sigma_{i} & =\sigma_{i} \sigma_{j+1}, \text { если } i \leqslant j ; \\

\sigma_{j} \delta_{i} & =\left\{\begin{array}{ll}

\delta_{i} \sigma_{j-1}, & \text { если } i<j, \\

\text { id, } & \text { если } i=j \text { или } i=j+1, \\

\delta_{i-1} \sigma_{j}, & \text { если } i>j+1 ;

\end{array}\right\}

\end{aligned}

\]



причем любое соотношение между этими отображениями является следствием соотношений (*). Это означает, что С. о. $X$ можно отожкдествить с системой $\left\{X_{n}, d_{i}, s_{i}\right\}$, состоящей из объектов $X_{n}, n \geqslant 0$, категории $\mathscr{C}$ и морфизмов $d_{i}: X_{n} \rightarrow X_{n-1}, 0 \leqslant i \leqslant n$, и $s_{i}: \quad X_{n} \rightarrow X_{n+1}, \quad 0 \leqslant i \leqslant n, \quad$ удовлетворяющих соотношениям



\[

d_{i} s_{j}=\left\{\begin{array}{c}

d_{i} d_{j}=d_{j-1} d_{i}, \text { если } i<j ; \\

s_{i-1} s_{j}=s_{j+1} s_{i}, \text { если } i \leqslant j ; \\

\text { id, если } i<j, \\

s_{j} d_{i-1}, \text { если } i>j \text { или } i=j+1,

\end{array}\right.

\]



Аналогично, косимплициальныӥ объект $X$ можно рассматривать как систему $\left\{X_{n}, d^{i}, s^{i}\right\}$, состоящую из объектов $X^{n}, n \geqslant 0\left(n-x\right.$ к о с л о е в), и морфизмов $d^{i}$ : $X^{n-1} \rightarrow X^{n}, 0 \leqslant i \leqslant n($ о п е р а т о р о в к о г р а н е й $)$, и $s^{i}: X^{n+1} \rightarrow X^{n}, 0 \leqslant i \leqslant n \quad($ п е р а т о р о в к ов ы р о ж д н и я), удовлетворяющих соотношениям (*) (в к-рых положено $\left.\delta_{i}=d^{i}, \sigma_{i}=s^{i}\right)$.



Сим плициальным отображен ием $f$ : $X \rightarrow Y$ С. о. $X$ в С. о. $Y$ (одной и той же категории $\mathscr{C})$ наз. произвольное преобразование (морфизм) функтора $X: \Delta \rightarrow \mathscr{C}$ в функтор $Y: \Delta \rightarrow \mathscr{C}$, т. е. такая система морфизмов $f_{n}: X_{n} \rightarrow Y_{n}, n \geqslant 0$, критерии $\mathscr{C}$, ЧTO



\[

\begin{gathered}

d_{i} f_{n+1}=f_{n} d_{i}, \quad 0 \leqslant i \leqslant n+1, \\

s_{i}^{\star} f_{n}=f_{n+1} s_{i}, \quad 0 \leqslant i \leqslant n .

\end{gathered}

\]



Симплициальные объекты категории $\mathscr{C}$ и их спмпліциальные отображения образуют категорию $\Delta^{\circ} \mathscr{C}$.



С им п ли ци а льн о й гомото п и е ї $h: f \simeq g$, связывающей симплициальные отображения $f, g: X \rightarrow Y$ симплициальных объектов категории $\mathrm{C}$, наз. такое семейство морфиизов $h_{i}: X_{n} \rightarrow Y_{n+1}, 0 \leqslant i \leqslant n$, категории $\mathscr{C}$, что



\[

\begin{gathered}

d_{0} h_{0}=f_{n} ; \\

d_{n} h_{n}=g_{n} ; \\

d_{i} h_{j}=\left\{\begin{array}{l}

h_{j-1} d_{i}, \text { если } i<j, \\

d_{j} h_{j-1}, \text { если } i=j>0, \\

h_{j} d_{i-1}, \text { если } i>j+1 ;

\end{array}\right. \\

s_{i} h_{j}=\left\{\begin{array}{lll}

h_{j+1} s_{i}, & \text { еслі } i \leqslant j, \\

h_{j} s_{i-1}, & \text { если } i>j .

\end{array}\right.

\end{gathered}

\]



На основе этого определения в категории $\Delta^{\circ} \mathscr{C}$ над произвольной категорией $\mathscr{C}$ можно воспроизвести по существу всю обычную теорию гомотопий. В случае категории множеств или топологич. пространств функтор геометрич. реализации (см. Силплициальное множество) переводит эту «симплициальную" теорию в обычную.



П р и м е р ы С. о.: симплициальное множество, симплициальное топологич. пространство, симплициальное алгебраич. многообразие, симилициальная группа, симплициальная абелева группа, симплициальная алгебра Ли, симплициальное гладкое многообразие и т. д. Каждая симплициальная абелева группа является цепным комплексом с граничным оператором $d=\Sigma(-1)^{i} d_{i}$.



Лum.: [1] $\Gamma$ а 6 р и е л ь П., Ц и с м а н М., Категории частных и теории гомотопий, пер. с англ., М., 1971; [2] M а у J. P., $\mathrm{m}$ ot $\mathrm{k}$ e K., Semisimpliziale algebraische Topologie, B.-



СИМПСОНА ФОРМУЛА - частный М. М. Постников. на- Котеса квадратурной форлулы, в к-рой берутся три узла:



\[

\int_{a}^{b} f(x) d x \cong \frac{b-a}{6}\left[f(a)+4 f\left(\frac{a+b}{2}\right)+f(b)\right]

\]



Пустъ промежуток $[a, b]$ разбит на $n$ частичных промежутков $\left[x_{i}, x_{i+1}\right], i=0,1,2, \ldots, n-1$, длины $h=$ $=(b-a) / n$, при этом $n$ считается четным числом, и для вычисления интеграла по промежутку $\left[x_{2 k}, x_{2 k+2}\right]$. исшользована квадратурная формула (1):



\[

\int_{x_{2 k}}^{x_{2 k+2}} f(x) d x \cong \frac{h}{3}\left[f\left(x_{2 k}\right)+4 f\left(x_{2 k+1}\right)+f\left(x_{2 k+2}\right)\right]

\]



Суммирование по $k$ от 0 до $n / 2-1$ левой и правой частей этого равенства приводит к составной С. ф.:



\[

\begin{aligned}

& \int_{a}^{b} f(x) d x \cong \frac{h}{3}\left\{f(a)+f(b)+2\left[f\left(x_{2}\right)+f\left(x_{4}\right)+\ldots\right.\right. \\

& \left.\left.\ldots+f\left(x_{n-2}\right)\right]+4\left[f\left(x_{1}\right)+f\left(x_{3}\right)+\ldots+f\left(x_{n-1}\right)\right]\right\},

\end{aligned}

\]



где $x_{j}=a+j h, j=0,1,2, \ldots, n$. Квадратурную формулу (2) также называют С. ф. (без добавления слова составная). Алгебраич. степень точности квадратурной формулы (2), как и формулы (1), равна 3.



Если подинтегральная функция $f(x)$ имеет непрерывную производную 4-го порядка на $[a, b]$, то погрешность $R(f)$ квадратурной формулы $(2)$ - разность между левой и правой частями приближенного равенства (2) - имеет представление



\[

R(f)=-\frac{b-a}{180} h^{4} f^{(\mathrm{IV})}(\xi),

\]



где $\xi$ - нек-рая точка из промежутка $[a, b]$.



C. ф. названа по имени Т. Симпсона (Th. Simpson), получившего ее в 1743, хотя эта формула была известна ранее, напр. Дж. Грегори (Ј. Gregory, 1668).



СИМСОНА ПРЯМАЯ - грямая, на к-рой лежат основания перпендикуля ров, опущенных из произвольной точки $P$ окружности, описанной вокруг треугольпика, на его стороны. С. п. делит на две равные части отрезок, соединяющий точку $P$ с точкой пересечения высот треугольника. С. п. названа по имени Р. Симсона (R. Simson), хотя она впервые была указана ранее.



СИМУЛА (от англ. «SIMUlation LAnguage», т. e. «язык моделирования») - название двух алгоритмич. языков, разработанных на основе алгола в Норвежском вычислительном центре и неофициально различаемых как симула 1 и симула-67.



С и м у л а 1 - проблемно-ориентированный язык для моделирования систем $\mathrm{c}$ дискретными событиями (напр., систем массового обслуживания), разработан в 1964. Спецификация модели сопоставляет компонентам системы (клиентам, станкам, материалам и т. п.) п р о ц е с с ы. Процесс имеет атрибуты (структуру данных) и программу действий (алгоритм). Модель работает по принципу к в а зи п а р а л ле ли з м а : в каждый момент активен только один іроцесс; исполняя свою программу, он может использовать свои и чужие атрибуты, порождать новые процессы, планировать себе и другим процессам с о б ы т и я- новые фазы активности (применяя встроенное в язык понятие дискретного времени), приостановить себя. Реализация С.1 шривела к разработке алгоритмич.





\end{document}